\documentclass[exercice]{thedocument}%
\usepackage{multicol}%
\startdatakeys
\source{}
\BO{2026}
\cycle{4}
\competencesDuSocle{CA}
\objectifsApprentissage{A1dd,E1ac,D2ac}
\automatismes{A1cg,A1ck}
\enddatakeys
\begin{document}%
\begin{flushleft}
    Un virus informatique infecte un fichier du disque dur s'un ordinateur
    et le rend inutilisable.\par La quantité totale de fichiers infectés
    décuple (c'est à dire est multipliée par 10) toutes les heures.
    Exprimer sous la forme d'une puissance de 10 le nombre total de fichier
    infectés au bout de :
\end{flushleft}
\begin{enumerate}
    \begin{multicols}{2}
        \item 3 heures
        \item cinq heures
        \item douze heures
        \item une journée
    \end{multicols}
\end{enumerate}
\end{document}
